\setAuthor{Valter Kiisk}
\setRound{piirkonnavoor}
\setYear{2019}
\setNumber{G 4}
\setDifficulty{3}
\setTopic{Termodünaamika}

\prob{Vaakum}
Vaatleme mõlemast otsast õhukindlalt suletud klaastoru pikkusega $\ell=\SI{1}{m}$ ning sisediameetriga $d=\SI{1}{cm}$. Õhurõhk on $p=\SI{101.3}{kPa}$.\\
\osa Kui palju tööd tuleb minimaalselt teha torus vaakumi tekitamiseks?\\
\osa Vakumeeritud toru hoitakse horisontaalselt ning selle sees ühte otsa on paigutatud teraskuul. Teraskuul saab torus libiseda hõõrdevabalt ning selle diameeter on võrdne toru sisediameetriga. Õnnetuse tõttu puruneb toru kuulikese poolne ots ning õhu surve paneb kuulikese liikuma vakumeeritud osa suunas. Kui suure kiiruse omandab kuulike vahetult enne teise otsa jõudmist? Terase tihedus on $\SI{7.9}{g/cm^3}$ ning võite eeldada, et $\ell \gg d$.\hint
Toru tühjakspumpamise käigus teeb torus sees olev õhk tööd välisrõhu $p$ vastu vahemaa $\ell$ vältel. Teises osas läheb sama hulk tööd teraskuuli kineetiliseks energiaks.\solu
\osa Toru tühjakspumpamine tähendab sisuliselt õhu väljasurumist torust, tehes tööd välisrõhu $p$ vastu. Õhk mõjub kuulikesele jõuga $pS$, kus toru ristlõikepindala $S=\pi(d/2)^2=\SI{0.79}{cm^2}$. Jõudu $pS$ tuleb rakendada teepikkusel $\ell$, nii et tehtud töö $A=pS\ell\approx\SI{8.0}{J}$. \\
\osa Kuulikese mass on
\[
m=\rho V=\rho\frac{4}{3}\pi\left(\frac{d}{2}\right)^3 = \SI{4.1}{g}.
\]
Kuna kuulike liigub eeldatavasti heli kiirusest hulga aeglasemalt, siis mõjub talle praktiliselt konstantne kiirendav jõud $pS$, jällegi teepikkusel $\ell$. Järelikult eelmises punktis leitud töö $A$ annab ühtlasi kuulikese kineetilise energia $mv^2/2$ toru teises otsas. Seega,
\[
v=\sqrt{\frac{2A}{m}} = \approx\SI{62}{m/s}.
\]
Alternatiivselt on lõppkiirus leitav üldtuntud valemi $v^2-v_0^2=2a\ell$ kaudu, kus algkiirus $v_0=0$ ja kiirendus $a=pS/m$.\probend